% AC-S52: scoped abstract commutator obstruction; no actual-Weil conclusion.
\section{The exact rank-one commutator does not supply a shrinking-gap rate}
\label{sec:spectral-parity-commutator-obstruction}

This is an abstract operator diagnosis, not a counterexample to the
actual localized Weil form.  It tests precisely whether compactness,
strong convergence of the normalized metric, strict compact metric
positivity, and the exact rank-one identity of AC-S51 force the
Hilbert--Schmidt comparison when the odd relative gap tends to zero.
They do not.  No claim is made that the example has the actual compact
spatial supports, Fourier-evaluation dictionary, or the arithmetic
explicit formula.

All Hilbert inner products below are linear in the first variable.
For a unit vector \(v\), write \(P_v f=\langle f,v\rangle v\).
Let \(\mathcal H_+=\mathcal H_-=\ell^2(\mathbb N)\) and set
\[
 D=\operatorname{diag}(2^{-j})_{j\ge1},\qquad
 C=\begin{pmatrix}0&D\\D&0\end{pmatrix}.
 \tag{PA.1}
\]
The letter \(D\) in this research example is the bounded off-diagonal
block, not the unbounded derivative denoted by \(D\) in AC-S51.
The fixed \(C\) is selfadjoint, injective, reflection-odd and
Hilbert--Schmidt, with \(\|C\|_{\mathcal S_2}^2=2/3\).
For \(n\ge1\), put
\[
 \tau=2^{-n},\quad \kappa=\tau^4,\quad r=1-\kappa,
 \qquad
 (v_n)_{n+k}=2^{-(k+1)/2}\ (k\ge0),
 \quad (v_n)_j=0\ (j<n).
 \tag{PA.2}
\]
Thus \(\|v_n\|=1\), and \(v_n\notin\operatorname{ran}D\).
Define the positive compact operator and positive scalars
\[
 H_n=\tau^4D^2,\qquad
 a_n=\langle (I-H_n)^{-1}v_n,v_n\rangle^{-1},\qquad
 b_n=\langle (rI-H_n)^{-1}Dv_n,Dv_n\rangle^{-1}.
 \tag{PA.3}
\]
The inverses exist since
\(\|H_n\|=\tau^4/4<r\).  Finally set
\[
 M_{n,+}=H_n+a_nP_{v_n},\qquad
 M_{n,-}=H_n+b_n(Dv_n)\otimes(Dv_n)^*,\qquad
 S_{n,\pm}=I-M_{n,\pm}.
 \tag{PA.4}
\]
Here the second rank-one term means
\(f\mapsto b_n\langle f,Dv_n\rangle Dv_n\).
Let \(e_n\) and \(f_n\) be the positive normalizations of
\[
 (I-H_n)^{-1}v_n,\qquad (rI-H_n)^{-1}Dv_n,
 \tag{PA.5}
\]
respectively.

\begin{proposition}[AC-S52: exact commutator without uniform comparison]
\label{AC-S52}
The operators in (PA.1)--(PA.5) have all the following properties.
Both \(M_{n,+}\) and \(M_{n,-}\) are compact and strictly positive
on every nonzero vector.  Their largest eigenvalues are simple and
equal to \(1\) and \(1-\kappa\), respectively.  They converge
strongly to zero.  In particular
\[
 \|M_{n,-}D\|_{\mathcal S_2}
       +\|DM_{n,+}\|_{\mathcal S_2}\longrightarrow0.
 \tag{PA.6}
\]
There is a bounded linear functional \(\ell_n\), zero on the odd
space and with \(\ell_n(e_n)=1\), such that the exact identity
\[
 DM_{n,+}-M_{n,-}D
       =S_{n,-}De_n\otimes\ell_n
 \tag{PA.7}
\]
holds.  The actual algebraic quotient of
\[
 K_n=C(I-e_n\otimes\ell_n)
 \tag{PA.8}
\]
by \(\mathbb C e_n\), with metric \(I-M_n\), is compact,
selfadjoint and injective.  In its isometric coordinates it is
\[
 B_n=\begin{pmatrix}
 0&S_{n,+}^{1/2}D S_{n,-}^{-1/2}\\
 S_{n,-}^{-1/2}D S_{n,+}^{1/2}&0
 \end{pmatrix},
 \tag{PA.9}
\]
restricted off the additional zero vector \(e_n\).
Nevertheless
\[
       \|B_n\|\ge \frac1{\sqrt5\,\tau}
                    \longrightarrow\infty.
 \tag{PA.10}
\]
Thus even operator-norm convergence to the fixed \(C\) fails,
although the inputs of the compact-product argument and the exact
commutator hold.
\end{proposition}

\begin{proof}
\proofstep{1. Positivity, top eigenvalues and strong convergence.}
For any positive \(A\), vector \(u\), and \(t>\|A\|\), the
operator \(A+c\,u\otimes u^*\), where
\(c=\langle(tI-A)^{-1}u,u\rangle^{-1}\), is bounded above by
\(tI\).  Indeed
\[
 tI-A-c\,u\otimes u^*
 =(tI-A)^{1/2}
 \left(I-c\,w\otimes w^*\right)(tI-A)^{1/2},
 \qquad w=(tI-A)^{-1/2}u,
\]
and \(c\|w\|^2=1\).  Its kernel is exactly the line
\(\mathbb C(tI-A)^{-1}u\).  Apply this with \((t,u)=(1,v_n)\)
and \((r,Dv_n)\).  This proves the assertions about the simple
largest eigenvalues and (PA.5).  Both metrics are compact because
\(H_n\) is compact and the perturbations have rank one.  They
are strictly positive because every diagonal entry of \(H_n\)
is positive.

The tail support gives \(v_n\rightharpoonup0\) and
\(Dv_n/\|Dv_n\|\rightharpoonup0\).  Also \(0<a_n\le1\) and
\(b_n\|Dv_n\|^2\le r\), directly from (PA.3).
Since \(\|H_n\|\to0\), (PA.4) implies strong convergence of
both \(M\) blocks to zero.  Finite-rank approximation of the fixed
Hilbert--Schmidt \(D\), followed by the uniform bounds
\(\|M_{n,\pm}\|\le1\), proves (PA.6).

\proofstep{2. Exact commutator and nonzero normalization.}
Write \(v=v_n\), \(a=a_n\), \(b=b_n\), and suppress \(n\)
in the following calculation.  Because \(H\) commutes with \(D\),
\[
 DM_+-M_-D=Dv\otimes\ell^0,
 \qquad
 \ell^0(h)=a\langle h,v\rangle-b\langle h,D^2v\rangle.
 \tag{PA.11}
\]
Evaluating this at \(e\), for which \(M_+e=e\), gives
\(S_-De=Dv\,\ell^0(e)\).  The left side is nonzero because
\(S_-\ge\kappa I>0\) and \(D\) is injective.  Hence
\(\ell^0(e)\ne0\).  Set
\(\ell=\ell^0/\ell^0(e)\) on the even space and zero on the
odd space.  It is bounded at each \(n\), and (PA.7) follows
with exactly the normalization used in AC-S51.

Equivalently
\[
 D(I-e\otimes\ell)=S_-^{-1}D S_+.
 \tag{PA.12}
\]
Thus \(K_n\) kills \(e\) and induces the stated quotient map.
On \(e^\perp\), \(S_+\) is boundedly invertible, since the
largest eigenvalue of the compact \(M_+\) is simple.  The
isometry with blocks \(S_+^{1/2}\) and \(S_-^{1/2}\)
therefore gives exactly (PA.9).  Its blocks are adjoints, so
it is compact selfadjoint.

For injectivity, note that \(e\notin\operatorname{ran}D\).
Its nonzero tail coordinates are a fixed scalar multiple of
\(v_j/(1-\tau^4 2^{-2j})\); dividing by \(2^{-j}\) gives a
sequence not in \(\ell^2\).
The lower-left block in (PA.9) has kernel exactly \(\mathbb C e\).
If its adjoint annihilates \(h\), then
\(D S_-^{-1/2}h\in\mathbb C e\).  Since \(e\notin\operatorname{ran}D\),
this vector is zero; injectivity of both factors gives \(h=0\).
This proves injectivity on the quotient.

\proofstep{3. Explicit amplification on the weakly escaping odd line.}
All vectors \(v,e,f\) are supported on the tail \(j\ge n\).
On that tail,
\[
       \|D\|=\tau,\qquad \|H\|=\eta:=\tau^6.
 \tag{PA.13}
\]
Put \(f^0=Dv/\|Dv\|\).  Three geometric series give
\[
 \|Dv\|^2=\frac47\tau^2,
 \qquad \|D^2v\|^2=\frac{16}{31}\tau^4,
 \qquad
 \langle Df^0,(I-P_v)Df^0\rangle
       =\frac{72}{217}\tau^2.
 \tag{PA.14}
\]
Since \(r\ge15/16\), normalization of
\((I-H/r)^{-1}f^0\) gives
\[
 \|f-f^0\|\le\frac{2(\eta/r)}{1-\eta/r}
                 \le3\eta .
 \tag{PA.15}
\]
For clarity, if \(w=(I-H/r)^{-1}f^0\), then
\(\|w-f^0\|\le(\eta/r)/(1-\eta/r)\) and
\(\|w/\|w\|-f^0\|\le2\|w-f^0\|\), proving the first
inequality.  The second follows from \(\eta\le1/64\).

The scalar \(a\) lies in \([1-\eta,1]\), so on the same tail
\[
       \|S_+-(I-P_v)\|\le2\eta.
 \tag{PA.16}
\]
Using \(\|Df\|,\|Df^0\|\le\tau\), equations
(PA.15)--(PA.16) imply
\[
 \left|\langle Df,S_+Df\rangle
       -\langle Df^0,(I-P_v)Df^0\rangle\right|
       \le 2\eta\tau^2+2\tau^2\|f-f^0\|
       \le8\tau^8.
 \tag{PA.17}
\]
Consequently, since \(\tau\le1/2\),
\[
 \langle Df,S_+Df\rangle
 \ge\left(\frac{72}{217}-\frac18\right)\tau^2
 =\frac{359}{1736}\tau^2
 >\frac15\tau^2.
 \tag{PA.18}
\]
But \(S_-f=\kappa f\), and hence the norm of (PA.9) on the
unit odd vector \(f\) is
\[
 \|B_n(0,f)\|^2
   =\kappa^{-1}\|S_+^{1/2}Df\|^2
   \ge\frac1{5\tau^2}.
\]
This proves (PA.10).
\end{proof}

\paragraph{The amount of primitive structure retained.}
The example also has exact metric symmetry on the mean-zero
hyperplane, not just a visually similar rank-one formula.
For \(M=M_+\oplus M_-\), equation (PA.11) and its adjoint give
\[
 \langle MC u,v\rangle=\langle M u,Cv\rangle
       \quad (\ell(u)=\ell(v)=0).
 \tag{PA.19}
\]
Indeed the lower-left block of \(CM-MC\) factors through
\(\ell\), and its upper-right block is minus the adjoint.
If desired, choose any \(\alpha_n>0\) tending to zero and use
\(J_n=\alpha_n^{-1/2}M_n^{1/2}\) as the abstract spatial
inclusion.  It is compact and injective.  The minima of
\(\|u\|^2/\|J_nu\|^2\) in the two parities are exactly
\(\alpha_n\) and \(\alpha_n/r\).  Equation (PA.19) is then the
corresponding integration-by-parts identity for the bounded
primitive \(C\) on \(\ker\ell\).

One can also retain the inverse identity on a dense test subspace without
postulating a spatial derivative.  The unbounded selfadjoint
inverse \(\mathcal D=C^{-1}\) satisfies \(C\mathcal D\phi=\phi\)
on its domain.  The subspace of that domain on which
\(\ell(\mathcal D\phi)=0\) is dense in \(\mathcal H\): on the
only nontrivial parity this functional is unbounded in the
Hilbert norm, since its representing vector before applying
\(D^{-1}\) is \(av-bD^2v\notin\operatorname{ran}D\).
A nonzero unbounded linear functional on a dense linear domain
has dense kernel.  To see this directly, failure of density
would yield a nonzero continuous functional annihilating that
kernel; on the domain it is a scalar multiple of the given
functional, making the latter bounded, a contradiction.
The density just proved is in \(\mathcal H\), not in the graph
norm of \(\mathcal D\); no graph-norm core claim is made.
This assertion does not identify the test subspace with actual
compact smooth functions on an interval.

\paragraph{Consequence for the research route.}
In this example the low \(S_-\) directions escape weakly to zero,
and \(D\) suppresses them in ordinary norm.  The exact positive
numerator (PA.18) still decays only as \(\tau^2\), whereas the
inverse-square-root factor has \(\kappa=\tau^4\).  Thus the
amplification wins.  No estimate omitting a rate can be deduced
from (PE.14), strong \(M_L\to0\), and compactness alone, even
with strict compact metric positivity and an injective quotient.
Any unconditional comparison for the actual Weil family must
use a further property of that family.  This diagnosis proves
neither failure nor truth of the actual relative parity gap,
the AC-S48 ratio, the original \(\Xi\) limit, or spatial
\(L^2\) convergence.
